% TOP_PROBLEM: {"problemId": "problem.morrison-kawamata-cone-conjecture", "canonicalTitle": "Morrison-Kawamata Cone Conjecture", "releaseRank": 240, "primaryDomain": "geometry_topology", "primaryDomainLabel": "Geometry & topology", "displayStatus": "open_with_solved_subcases", "releaseStatus": "open_with_solved_subcases"}
% !TeX program = lualatex
% Definition and sources only; this file is not an LLM solution attempt.
\documentclass[11pt]{article}
\usepackage[margin=25mm]{geometry}
\usepackage{fontspec}
\setmainfont{Latin Modern Roman}
\usepackage{amsmath,amssymb,mathtools}
\usepackage{unicode-math}
\setmathfont{Latin Modern Math}
\usepackage{xurl,hyperref}
\hypersetup{colorlinks=true,urlcolor=blue}
\setcounter{secnumdepth}{0}
\setlength{\parindent}{0pt}
\setlength{\parskip}{5pt}
\setlength{\emergencystretch}{3em}
\begin{document}
\section*{\texorpdfstring{Morrison-Kawamata Cone Conjecture}{Morrison-Kawamata Cone Conjecture}}\label{240}
\subsection{Definitions and mathematical statement}
Let $(X,\Delta)$ be a projective complex ${\mathbb{Q}}$-factorial klt pair with $K_X+\Delta\equiv0$. Here ${\mathbb{Q}}$-factorial means every Weil divisor has a Cartier multiple, and klt means discrepancies exceed $-1$. In $N^1(X)_{{\mathbb{R}}}$, let $\operatorname{Nef}^e$ be the effective nef cone and $\operatorname{Mov}^e$ the effective movable cone in the cited convention. Automorphisms preserve the pair; pseudo-automorphisms are birational self-maps preserving it and isomorphic in codimension one. The two assertions ask that each relevant cone be covered by translates of a rational polyhedral cone under its indicated group, with interiors of distinct translates disjoint in the standard fundamental-domain sense. Rational polyhedral means generated by finitely many rational divisor classes. Both the nef and movable assertions are part of the target.
\par\smallskip\textit{Formulation and concept references:} \href{https://arxiv.org/abs/2406.07307}{[S1]}.

\subsection{Short English statement}
Do automorphisms and pseudo-automorphisms reduce the effective nef and movable cones of every klt Calabi–Yau pair to rational polyhedral fundamental domains?
\subsection{Sources}
\begin{itemize}
\item[S1]The effective cone conjecture for Calabi-Yau pairs.
\url{https://arxiv.org/abs/2406.07307}.
\textit{Formulation source.}
\end{itemize}
\end{document}
